\section{Unified Mobile Private Contacts and Recipient Resolution}
\label{sec:pilot-unified-mobile-contacts}

Pilot's unified mobile application now presents the existing Personal Pilot contact authority as
a dedicated private Contacts surface.  The mother brain remains the canonical owner of contacts
and routing information.  The phone provides possession-bound inspection and editing; shared
television surfaces receive neither email addresses nor telephone numbers and cannot silently
choose a recipient.

\subsection{Contact Model and Routes}

A contact belongs to exactly one private persona and may contain one or more independently
validated routes:
\begin{itemize}
  \item a Pilot persona and mother-brain identity for encrypted Pilot-to-Pilot delivery;
  \item an international telephone number for an authorized WhatsApp handoff; and
  \item an email address for an authorized email adapter.
\end{itemize}

The owner may select a preferred route.  If the selected route is unavailable, Pilot does not
invent an endpoint.  Contact creation requires at least one usable route, validates email and
international-number syntax, and validates the shape of remote Pilot identifiers.  Contact data
is not itself authority to send: communication, task delegation and external-provider execution
retain their own proposal, approval and receipt boundaries.

\subsection{Mobile Experience}

The Personal mode now includes a Contacts shortcut.  A connected user can list and search their
own contacts, add a contact manually, edit an existing record, choose the preferred route, and
remove a contact.  Where the browser exposes the privacy-preserving Contact Picker API, the user
may deliberately select one phone contact and copy its proposed name, email and telephone number
into the form.  Pilot still requires the user to review and save that form.  Browsers without the
picker show a clear manual-entry path rather than requesting broad address-book access.

Pilot-to-Pilot identifiers are placed in a secondary disclosure panel because they are normally
learned through a trusted invitation rather than typed by ordinary users.  No contact operation
sends a message, creates a task or grants access to another mother brain.

\subsection{Possession and Privacy Boundary}

Contact operations use three separately leased capabilities:
\begin{itemize}
  \item \texttt{contacts.read} for listing private contact routes;
  \item \texttt{contacts.write} for creation, editing and removal; and
  \item \texttt{contacts.resolve} for turning a typed or recognized name into a proposed recipient.
\end{itemize}

Every request is signed by the phone's non-exportable Ed25519 key and validated against its active
certificate, persona and short-lived lease.  A phone cannot substitute another persona identifier.
The general Personal Pilot summary exposes only contact names and available route types; exact
addresses are returned only through the dedicated possession-authorized Contacts endpoint.

\subsection{Safe Recipient Resolution}

Resolution first attempts a normalized exact-name match.  If none exists, Pilot applies the
existing bounded phonetic matcher intended for local speech-recognition errors.  A unique phonetic
match may be proposed, but the result is explicitly labelled as phonetic.  Zero matches return
\texttt{not\_found}; multiple matches return \texttt{ambiguous} with a private choice list.  Pilot
never chooses the first ambiguous person.

The resolved record identifies the proposed route and recipient reference, but its
\texttt{external\_effect} value remains false.  A later communication or delegation stage must show the exact
recipient and content and obtain its own private approval before delivery.

\subsection{Durability, Removal and Historical Evidence}

Mobile contact writes use idempotency keys bound to a canonical hash of the exact owner, contact
and route fields.  An identical retry returns the original contact.  Reusing the key with changed
content fails closed.  Removal is a soft revocation: the contact disappears from active search and
resolution while historical tasks, messages and delivery receipts continue to reference the
original contact identifier.  Repeated removal returns the existing result instead of mutating
history twice.

\subsection{Verification}

Automated evidence covers persona filtering, exact-route visibility on the private phone, contact
creation, repeat-request deduplication, unique phonetic resolution, unavailable-route rejection,
soft removal, removal replay and cross-person denial.  Surface checks verify the dedicated
Contacts experience and all signed mother-brain routes.  At coded closure of \texttt{MOB-0505},
all 458 AION Fabric tests passed and the optimized Next.js application compiled and statically
generated the unified mobile route.  Real WhatsApp, email and Pilot-to-Pilot delivery remain
separate governed execution stages; this capability resolves recipients but deliberately sends
nothing.
